% GENERATED FILE. Edit canonical lessons, then run npm run generate:books.
% canonical-source-hash: fnv1a64:1a2f7a23003d91cb
% canonical-lessons: JA-C219-korobu, JA-C219-suberu, JA-C219-oreru, JA-C219-taoreru, JA-C219-kega

\chapter{Slips, Falls and Breaks}
\label{ch:ja-slips-falls-and-breaks}
\hlchaptermodality{219}

\begin{chapteropening}
\textbf{By the end of this chapter:} \emph{I can ask whether someone fell or got hurt.}

Complete the last lesson of chapter 219: I can ask whether someone fell or got hurt.
\end{chapteropening}

\section[\texorpdfstring{\ja{ころぶ} (korobu)}{ころぶ (korobu)}]{\ja{ころぶ} (korobu) — to fall over}

\label{lesson:JA-C219-korobu}



\begin{quote}
Before the new one: say the Japanese for a wound, then the Japanese for life.
\end{quote}

\subsection*{You'll want to know: \ja{ころぶ}}
\textbf{\ja{ころぶ}} — \emph{korobu} — \textquotedblleft{}to fall over\textquotedblright{}.

\begin{quote}
\textbf{\ja{こどもは} \ja{ころびませんでしたか。}} — \emph{kodomo wa korobimasen deshita ka} — \textquotedblleft{}Didn't the child fall over?\textquotedblright{}
\end{quote}

Polite \textbf{\ja{ころびます}}, \emph{korobimasu}; a godan verb.

\begin{cousinweb}
\textbf{\ja{ころぶ}} (\emph{korobu}) is a native Japanese word (wago), inherited rather than borrowed from Chinese; its kanji was attached later as a kun'yomi, and its older origin is not settled (JMdict, KANJIDIC2).
\end{cousinweb}

\begin{grammarlens}[title={What you've built}]
One more word for everyday life at A2.
\end{grammarlens}

\subsection*{Your turn}
\begin{itemize}
\raggedright
  \item \emph{Say it:} \emph{korobu}
  \item \emph{Say it:} \emph{korobu}, once more
  \item \emph{Say it:} \emph{kodomo wa korobimasen deshita ka} — Didn't the child fall over?
  \item \emph{Recall:} say the Japanese for a wound, then the Japanese for life, then say \emph{korobu} again
\end{itemize}

\begin{tcolorbox}[breakable,colback=teal!4,colframe=teal!35!black,title={Before you move on}]
What does \ja{ころぶ} mean? (\textquotedblleft{}To fall over\textquotedblright{}.) Say it once more.
\end{tcolorbox}
\section[\texorpdfstring{\ja{すべる} (suberu)}{すべる (suberu)}]{\ja{すべる} (suberu) — to slip; to slide}

\label{lesson:JA-C219-suberu}



\begin{quote}
Before the new one: say the Japanese for life, then the Japanese for to fall over.
\end{quote}

\subsection*{You'll want to know: \ja{すべる}}
\textbf{\ja{すべる}} — \emph{suberu} — \textquotedblleft{}to slip; to slide\textquotedblright{}.

\begin{quote}
\textbf{\ja{ゆきの} \ja{みちは} \ja{すべりますか。}} — \emph{yuki no michi wa suberimasu ka} — \textquotedblleft{}Is the snowy road slippery?\textquotedblright{}
\end{quote}

Polite \textbf{\ja{すべります}}, \emph{suberimasu}; a godan verb.

\begin{cousinweb}
\textbf{\ja{すべる}} (\emph{suberu}) is a native Japanese word (wago), inherited rather than borrowed from Chinese; its kanji was attached later as a kun'yomi, and its older origin is not settled (JMdict, KANJIDIC2).
\end{cousinweb}

\begin{grammarlens}[title={What you've built}]
One more word for everyday life at A2.
\end{grammarlens}

\subsection*{Your turn}
\begin{itemize}
\raggedright
  \item \emph{Say it:} \emph{suberu}
  \item \emph{Say it:} \emph{suberu}, once more
  \item \emph{Say it:} \emph{yuki no michi wa suberimasu ka} — Is the snowy road slippery?
  \item \emph{Recall:} say the Japanese for life, then the Japanese for to fall over, then say \emph{suberu} again
\end{itemize}

\begin{tcolorbox}[breakable,colback=teal!4,colframe=teal!35!black,title={Before you move on}]
What does \ja{すべる} mean? (\textquotedblleft{}To slip; to slide\textquotedblright{}.) Say it once more.
\end{tcolorbox}
\section[\texorpdfstring{\ja{おれる} (oreru)}{おれる (oreru)}]{\ja{おれる} (oreru) — to snap; to break (of a bone or stick)}

\label{lesson:JA-C219-oreru}



\begin{quote}
Before the new one: say the Japanese for to fall over, then the Japanese for to slip.
\end{quote}

\subsection*{You'll want to know: \ja{おれる}}
\textbf{\ja{おれる}} — \emph{oreru} — \textquotedblleft{}to snap; to break (of a bone or stick)\textquotedblright{}.

\begin{quote}
\textbf{\ja{ほねが} \ja{おれましたか。}} — \emph{hone ga oremashita ka} — \textquotedblleft{}Did you break a bone?\textquotedblright{}
\end{quote}

Polite \textbf{\ja{おれます}}, \emph{oremasu}; an ichidan verb.

\begin{cousinweb}
\textbf{\ja{おれる}} (\emph{oreru}) is a native Japanese word (wago), inherited rather than borrowed from Chinese; its kanji was attached later as a kun'yomi, and its older origin is not settled (JMdict, KANJIDIC2).
\end{cousinweb}

\begin{grammarlens}[title={What you've built}]
One more word for everyday life at A2.
\end{grammarlens}

\subsection*{Your turn}
\begin{itemize}
\raggedright
  \item \emph{Say it:} \emph{oreru}
  \item \emph{Say it:} \emph{oreru}, once more
  \item \emph{Say it:} \emph{hone ga oremashita ka} — Did you break a bone?
  \item \emph{Recall:} say the Japanese for to fall over, then the Japanese for to slip, then say \emph{oreru} again
\end{itemize}

\begin{tcolorbox}[breakable,colback=teal!4,colframe=teal!35!black,title={Before you move on}]
What does \ja{おれる} mean? (\textquotedblleft{}To snap; to break (of a bone or stick)\textquotedblright{}.) Say it once more.
\end{tcolorbox}
\section[\texorpdfstring{\ja{たおれる} (taoreru)}{たおれる (taoreru)}]{\ja{たおれる} (taoreru) — to collapse; to fall down}

\label{lesson:JA-C219-taoreru}



\begin{quote}
Before the new one: say the Japanese for to slip, then the Japanese for to snap.
\end{quote}

\subsection*{You'll want to know: \ja{たおれる}}
\textbf{\ja{たおれる}} — \emph{taoreru} — \textquotedblleft{}to collapse; to fall down\textquotedblright{}.

\begin{quote}
\textbf{\ja{かぜで} \ja{きが} \ja{たおれました。}} — \emph{kaze de ki ga taoremashita} — \textquotedblleft{}A tree fell down in the wind.\textquotedblright{}
\end{quote}

Polite \textbf{\ja{たおれます}}, \emph{taoremasu}; an ichidan verb. Its partner \textbf{\ja{たおす}}, \emph{taosu}, means \textquotedblleft{}to knock down\textquotedblright{}.

\begin{cousinweb}
\textbf{\ja{たおれる}} (\emph{taoreru}) is a native Japanese word (wago), inherited rather than borrowed from Chinese; its kanji was attached later as a kun'yomi, and it pairs with \emph{taosu}, \textquotedblleft{}to knock down\textquotedblright{} (JMdict, KANJIDIC2).
\end{cousinweb}

\begin{grammarlens}[title={What you've built}]
One more word for everyday life at A2.
\end{grammarlens}

\subsection*{Your turn}
\begin{itemize}
\raggedright
  \item \emph{Say it:} \emph{taoreru}
  \item \emph{Say it:} \emph{taoreru}, once more
  \item \emph{Say it:} \emph{kaze de ki ga taoremashita} — A tree fell down in the wind.
  \item \emph{Recall:} say the Japanese for to slip, then the Japanese for to snap, then say \emph{taoreru} again
\end{itemize}

\begin{tcolorbox}[breakable,colback=teal!4,colframe=teal!35!black,title={Before you move on}]
What does \ja{たおれる} mean? (\textquotedblleft{}To collapse; to fall down\textquotedblright{}.) Say it once more.
\end{tcolorbox}
\section[\texorpdfstring{\ja{けが} (kega)}{けが (kega)}]{\ja{けが} (kega) — an injury}

\label{lesson:JA-C219-kega}



\begin{quote}
Before the new one: say the Japanese for to snap, then the Japanese for to collapse.
\end{quote}

\subsection*{You'll want to know: \ja{けが}}
\textbf{\ja{けが}} — \emph{kega} — \textquotedblleft{}an injury\textquotedblright{}.

\begin{quote}
\textbf{\ja{けがは} \ja{ありませんか。}} — \emph{kega wa arimasen ka} — \textquotedblleft{}Are you hurt?\textquotedblright{}
\end{quote}

\begin{cousinweb}
\textbf{\ja{けが}} (\emph{kega}) is a Japanese word usually written in kana (JMdict); whether it is native and how its kanji spelling arose are not settled.
\end{cousinweb}

\begin{grammarlens}[title={What you've built}]
One more word for everyday life at A2.
\end{grammarlens}

\subsection*{Your turn}
\begin{itemize}
\raggedright
  \item \emph{Say it:} \emph{kega}
  \item \emph{Say it:} \emph{kega}, once more
  \item \emph{Say it:} \emph{kega wa arimasen ka} — Are you hurt?
  \item \emph{Recall:} say the Japanese for to snap, then the Japanese for to collapse, then say \emph{kega} again
\end{itemize}

\begin{tcolorbox}[breakable,colback=teal!4,colframe=teal!35!black,title={Before you move on}]
What does \ja{けが} mean? (\textquotedblleft{}An injury\textquotedblright{}.) Say it once more.
\end{tcolorbox}
